% Options for packages loaded elsewhere
\PassOptionsToPackage{unicode}{hyperref}
\PassOptionsToPackage{hyphens}{url}
\documentclass[
]{article}
\usepackage{xcolor}
\usepackage[margin=1in]{geometry}
\usepackage{amsmath,amssymb}
\setcounter{secnumdepth}{-\maxdimen} % remove section numbering
\usepackage{iftex}
\ifPDFTeX
  \usepackage[T1]{fontenc}
  \usepackage[utf8]{inputenc}
  \usepackage{textcomp} % provide euro and other symbols
\else % if luatex or xetex
  \usepackage{unicode-math} % this also loads fontspec
  \defaultfontfeatures{Scale=MatchLowercase}
  \defaultfontfeatures[\rmfamily]{Ligatures=TeX,Scale=1}
\fi
\usepackage{lmodern}
\ifPDFTeX\else
  % xetex/luatex font selection
\fi
% Use upquote if available, for straight quotes in verbatim environments
\IfFileExists{upquote.sty}{\usepackage{upquote}}{}
\IfFileExists{microtype.sty}{% use microtype if available
  \usepackage[]{microtype}
  \UseMicrotypeSet[protrusion]{basicmath} % disable protrusion for tt fonts
}{}
\makeatletter
\@ifundefined{KOMAClassName}{% if non-KOMA class
  \IfFileExists{parskip.sty}{%
    \usepackage{parskip}
  }{% else
    \setlength{\parindent}{0pt}
    \setlength{\parskip}{6pt plus 2pt minus 1pt}}
}{% if KOMA class
  \KOMAoptions{parskip=half}}
\makeatother
\usepackage{graphicx}
\makeatletter
\newsavebox\pandoc@box
\newcommand*\pandocbounded[1]{% scales image to fit in text height/width
  \sbox\pandoc@box{#1}%
  \Gscale@div\@tempa{\textheight}{\dimexpr\ht\pandoc@box+\dp\pandoc@box\relax}%
  \Gscale@div\@tempb{\linewidth}{\wd\pandoc@box}%
  \ifdim\@tempb\p@<\@tempa\p@\let\@tempa\@tempb\fi% select the smaller of both
  \ifdim\@tempa\p@<\p@\scalebox{\@tempa}{\usebox\pandoc@box}%
  \else\usebox{\pandoc@box}%
  \fi%
}
% Set default figure placement to htbp
\def\fps@figure{htbp}
\makeatother
\setlength{\emergencystretch}{3em} % prevent overfull lines
\providecommand{\tightlist}{%
  \setlength{\itemsep}{0pt}\setlength{\parskip}{0pt}}
\usepackage{bookmark}
\IfFileExists{xurl.sty}{\usepackage{xurl}}{} % add URL line breaks if available
\urlstyle{same}
\hypersetup{
  hidelinks,
  pdfcreator={LaTeX via pandoc}}

\author{}
\date{\vspace{-2.5em}}

\begin{document}

FINAL R NOTEBOOK WITH CODE Karolina Szczepanik

\texttt{r\ library(dplyr)\ library(tidyverse)\ library(estimatr)\ library(Rcpp)\ library(stargazer)\ library(ggplot2)\ library(gridExtra)}r

\section{import movie dataset into R}\label{import-movie-dataset-into-r}

movies \textless- read.csv(``movies\_final.csv'')

\section{graph 1:compares awards and IMDb
ratings}\label{graph-1compares-awards-and-imdb-ratings}

g1 \textless- ggplot(movies, aes(Awards\_Count, imdbRating)) +
geom\_point() + labs(title = ``Awards vs IMDb Rating'', x = ``Awards'',
y = ``IMDb Rating'')

\section{graph 2: compares runtime and IMDb
ratings}\label{graph-2-compares-runtime-and-imdb-ratings}

g2 \textless- ggplot(movies, aes(Runtime, imdbRating)) + geom\_point() +
labs(title = ``Runtime vs IMDb Rating'', x = ``Runtime'', y = ``IMDb
Rating'')

\#puts both graphs together grid.arrange(g1, g2, ncol = 1)

\section{graph 3: shows where most IMDb ratings fall in the
dataset}\label{graph-3-shows-where-most-imdb-ratings-fall-in-the-dataset}

ggplot(movies, aes(imdbRating)) + geom\_density(fill = ``dark grey'',
color = ``black'', alpha = 0.7) + labs( title = ``Distribution of IMDb
Ratings'', x = ``IMDb Rating'', y = ``Frequency'' )

\section{regression 1: full dataset with more than 3
variables}\label{regression-1-full-dataset-with-more-than-3-variables}

model1 \textless- lm(imdbRating \textasciitilde{} Runtime +
Awards\_Count + Genre\_Count, data = movies) summary(model1)

\section{regression 2: filtered dataset, only movies longer than 120
minutes}\label{regression-2-filtered-dataset-only-movies-longer-than-120-minutes}

long\_movies \textless- movies{[}movies\$Runtime \textgreater{} 120, {]}

model2 \textless- lm(imdbRating \textasciitilde{} Awards\_Count +
Genre\_Count, data = long\_movies) summary(model2)

\section{regression 3: preferred
model}\label{regression-3-preferred-model}

model3 \textless- lm(imdbRating \textasciitilde{} Runtime +
Awards\_Count + Genre\_Count, data = movies) summary(model3)

Regression Interpretation:

Regression 1 --\textgreater{} looks at the relationship between IMDb
ratings, runtime, awards count, and genre count across the full movie
dataset. Runtime had a small positive relationship with IMDb ratings,
meaning longer movies were slightly associated with higher ratings.

Regression 2 --\textgreater{} only focuses on movies longer than 120
minutes. This was done to see if the relationships changed for longer
films. In this model, awards count and genre count did not show a strong
relationship with IMDb ratings.

Regression 3 --\textgreater{} This is my preferred model because it uses
the full dataset and includes multiple movie characteristics together.
Overall, the results suggest that runtime may have a stronger
relationship with IMDb ratings than awards count or genre count in this
dataset.

These regressions show correlations between variables, not causation.

\end{document}
